\chapter{人声编排：让主唱站在前面，让伴唱补在后面}
\label{ch:vocals}

\classquote{诗者，志之所之也。在心为志，发言为诗。}{《毛诗序》}{诗是心意所在：在心里叫志，说出来就成了诗。人声编曲的第一件事，是别让这句话被配角盖住。人声编曲的第一原则不是“加多少层”，而是“歌词还听不听得到”。主唱是主角，双唱、和声、叠唱、装饰句都是配角；配角抢戏，整段就会糊。}

\section{人声不是一条轨，是一组角色}
\label{sec:vocal-roles}

一首流行歌的人声通常包含这些角色：

\begin{itemize}
\item \textbf{主唱}：唱主歌词，承担叙事与情绪；它的音区、咬字和呼吸必须最清楚。
\item \textbf{双唱}：同一句歌词再唱一遍，叠在主唱上增加厚度；不是和声，唱的是同一条旋律。
\item \textbf{和声}：唱不同的音，通常在主唱下方三度/六度，或在副歌做高低两条线。
\item \textbf{叠唱}：同一句歌词叠三到五层，做成“人声墙”；常见于副歌最后一遍或尾奏。
\item \textbf{装饰句}：主唱句尾的短促回应、喊唱、气声；它不承担歌词信息，只加颜色。
\end{itemize}

这些角色可以出现在不同段落：主歌只要主唱，预副歌加一层双唱，副歌开和声，最后一遍副歌才叠唱。\textbf{不是每段都要全员上场。}

\section{主唱：先保证歌词清楚}
\label{sec:lead-vocal}

主唱编排的第一件事是“留白”：

\begin{itemize}
\item \textbf{句长}：一个乐句不超过四小节；超过就拆成两句，中间留半拍到一拍呼吸。
\item \textbf{音区}：主唱停在说话音高附近最自然；副歌可以抬高，但不要整首都在最高音区。
\item \textbf{咬字}：辅音要留时间，不要把三个字挤在一拍里；歌词密的地方，旋律要简化。
\item \textbf{空间}：主唱句尾留空，让和声/装饰句有地方进；主唱唱满，配角只能抢。
\end{itemize}

如果主唱轨已经唱得很满，正确的做法不是加压缩，而是删掉一些音符/让出句尾。人声编排首先是“让主唱能呼吸”。

\section{双唱：同一条旋律，第二次唱}
\label{sec:double}

\term{双唱}{double}是同一句歌词唱两遍，叠在一起。它的作用不是和声，而是让主唱更厚、更稳：

\begin{itemize}
\item \textbf{层数}：主歌不加或只加一层；副歌加一层；最后一遍副歌可以两层。
\item \textbf{录音}：真正的双唱要两次录音，不是复制同一条音频；复制同一条会变成“更响”，而不是“更厚”，还可能有相位问题。
\item \textbf{对齐}：两次录音不必逐字完全重合；让辅音稍微错开几毫秒，会更有“人”的感觉；错太多会变成合唱/回声。
\item \textbf{音区}：双唱唱同一旋律，不唱低八度；低八度属于和声/叠唱，不属于双唱。
\end{itemize}

如果只有一条主唱录音，宁可不加双唱，也不要用复制轨硬充厚度；复制轨只适合做“宽化”效果，不适合冒充两次演唱。

\section{和声：在主唱下方写第二条线}
\label{sec:harmony}

\term{和声}{vocal-harmony}与器乐和声不同：它是一条能被听出来的旋律线，不是把和弦音随便塞进人声轨。

\begin{itemize}
\item \textbf{音程}：最常用的是主唱下方三度或六度；五度会显得空，二度会打架；不确定时先用三度。
\item \textbf{音区}：和声通常在主唱下方，不要和主唱同音区抢；高八度只用在副歌尾句。
\item \textbf{走向}：和声线要平滑，不要每个字都跳；主唱长音时，和声可以稍微移动，制造张力。
\item \textbf{段落}：主歌不加和声，预副歌加一条低和声，副歌加两条（高三度/低三度）；桥段可以只留一条。
\item \textbf{歌词}：和声唱同样的歌词，或只唱元音；不要在和声里写主唱没有的新词，否则听众会分心。
\end{itemize}

\section{叠唱与装饰句：厚度和颜色}
\label{sec:stack-adlib}

\term{叠唱}{vocal-stacking}是把同一句歌词叠多层，做成“人声墙”。它适合副歌最后一遍、尾奏或需要仪式感的段落：

\begin{itemize}
\item \textbf{层数}：三到五层足够；超过五层，歌词会糊，且每层都变薄。
\item \textbf{分工}：两层唱主旋律，一层低三度，一层高三度，一层只唱元音/气声；不要五层都唱同一句。
\item \textbf{声像}：主唱居中，叠唱向左右铺开；不要所有层都挤在中间。
\item \textbf{段落}：叠唱只在最后一遍副歌或尾奏出现；从主歌就开始叠，到副歌就没有“更多”的感觉。
\end{itemize}

\term{装饰句}{ad-lib}是主唱句尾的短回应：一个“哦”、一声叹气、一句喊唱。它只占半拍到两拍，落在主唱留白处；装饰句太多，会变成第二个主唱。

\begin{figure}[htbp]
  \centering
  \begin{tikzpicture}[font=\small]
    \draw[->] (0,0) -- (9.2,0) node[right] {bars};
    \foreach \x in {1,...,8} {\draw[gray!30] (\x,0.1) -- (\x,3.6); \node[font=\scriptsize] at (\x,0.05) {\x};}
    \fill[yebanblue!40] (0.2,2.9) rectangle (8.2,3.25);
    \node[font=\scriptsize] at (4.2,3.07) {lead};
    \fill[yebanblue!65] (4.2,2.15) rectangle (8.2,2.5);
    \node[font=\scriptsize] at (6.2,2.32) {double};
    \fill[green!45] (4.2,1.4) rectangle (8.2,1.75);
    \node[font=\scriptsize] at (6.2,1.57) {harmony};
    \fill[orange!55] (6.2,0.65) rectangle (8.2,1.0);
    \node[font=\scriptsize] at (7.2,0.82) {stack};
    \node[font=\scriptsize, anchor=west] at (8.4,2.32) {chorus};
  \end{tikzpicture}
  \caption{八小节人声层进出：主唱全程，双唱与和声从副歌进入，叠唱只在最后两小节；主歌留白。}
  \label{fig:vocal-stack}
\end{figure}

\section{夜半落地：音频轨、片段与分轨}
\label{sec:vocal-yeban}

人声是音频，不是 MIDI；它在夜半里的编排方式与器乐不同：

\begin{enumerate}
\item \textbf{建音频轨}：用「轨道 → + 新建音频轨」建主唱、双唱、和声、叠唱各一条；不要把所有层塞进一条音频轨。
\item \textbf{录音或导入}：用「文件 → 实时录音至 WAV...」录主唱，或用「文件 → 导入音频文件...」导入已有录音；导入的素材先进素材池，再摆到音频轨。
\item \textbf{剪片段}：选中音频片段后拖边缘改长度，拖动改位置；主唱句尾留空，不要把整段剪得严丝合缝。
\item \textbf{复制与对齐}：双唱/和声如果是另一条录音，复制片段后对齐主唱；不要用同一条主唱复制来冒充双唱；对齐时允许辅音有几毫秒错开。
\item \textbf{用片段静音做编排}：不需要的层用片段静音开关关掉，而不是删掉；这样副歌最后一遍可以快速打开叠唱。
\item \textbf{用推子与声像摆位}：主唱居中，双唱/和声稍向左右，叠唱铺宽；混音台（X / F8）里逐层检查，不要一次全开。
\item \textbf{导出分轨}：用「文件 → 导出分轨（Stems）...」把人声各层单独导出，交给混音/母带；导出回执会报条数与跳过条数。
\end{enumerate}

\screenshot{arrangement-vocal-stack}{主唱、双唱、和声、叠唱四条音频轨同屏，片段进出与声像位置可见}

\section{常见错误与修正}
\label{sec:vocal-mistakes}

\begin{longtable}{p{0.34\textwidth}p{0.58\textwidth}}
\toprule
常见错误 & 修正 \\
\midrule
\endfirsthead
\toprule
常见错误 & 修正 \\
\midrule
\endhead
从主歌到副歌一直叠唱 & 主歌只留主唱；双唱/和声从副歌进，叠唱只在最后一遍。 \\
复制主唱冒充双唱 & 复制轨只做宽化；真正的双唱要两次录音。 \\
和声与主唱同音区 & 和声放到主唱下方三度/六度；高八度只用在尾句。 \\
和声每个字都跳 & 和声线要平滑；长音处才做移动。 \\
装饰句盖住主唱 & 装饰句只落主唱留白，半拍到两拍；不要连续回应。 \\
五层叠唱全唱主旋律 & 分层：主旋律、低三度、高三度、元音/气声；不要五层同词。 \\
人声全挤在中间 & 主唱居中，双唱/和声/叠唱向左右铺开。 \\
\bottomrule
\end{longtable}

\section{练习}
\label{sec:vocal-exercise}

\begin{exercise}
用八小节做一段“主唱＋双唱＋和声＋叠唱”的人声编排：
\begin{enumerate}
\item 主唱写四小节一句，句尾留一拍；第二句重复，但尾音改低。
\item 副歌第 5 小节加一层双唱（同旋律）；第 5--8 小节加一条低三度和声。
\item 最后两小节加两层叠唱：一层主旋律、一层低三度；不要五层同词。
\item 用片段静音开关做主歌/副歌对比；导出分轨，确认每层都能单独找到。
\end{enumerate}
最后 Solo 主唱：歌词是否每个字都清楚？如果和声/叠唱拿掉后歌词更清楚，说明它们已经完成了配角的任务。
\end{exercise}

\section{本章要点}
\label{sec:vocal-takeaways}

\begin{itemize}
\item 人声是一组角色：主唱、双唱、和声、叠唱、装饰句；不是一条轨加效果。
\item 主唱优先：句长、音区、咬字、留白，先保证歌词清楚。
\item 双唱是同旋律两次录音；复制同一条不是双唱。
\item 和声是第二条旋律线，通常在主唱下方三度/六度，走向平滑。
\item 叠唱三到五层，分工唱；只在最后副歌/尾奏出现。
\item 夜半里用多条音频轨、片段静音、推子/声像与分轨导出落地。
\end{itemize}

\section{扩展阅读}
\label{sec:09-vocals-further-reading}
\begin{itemize}
  \item \href{https://en.wikipedia.org/wiki/Vocal_harmony}{Vocal harmony（人声和声）}：看清和声与旋律的主从关系。
  \item \href{https://en.wikipedia.org/wiki/Mixing_(music)}{Mixing（混音）}：了解频段让位的基本思路。
  \item \href{https://www.youtube.com/watch?v=ADPbHK3Wi4k}{一条人声做出多层和声}：看叠唱与和声层怎样搭起来。
  \item \href{https://www.youtube.com/watch?v=DSc9i4DpjQE}{主唱与和声的频段分工}：用总线处理给主唱留出中频位置。
\end{itemize}

\begin{helpmapping}
\begin{longtable}{p{0.45\textwidth}p{0.45\textwidth}}
\toprule
小节 & 帮助锚点 \\
\midrule
人声不是一条轨，是一组角色 & \texttt{arrangement-vocal-roles} \\
主唱：先保证歌词清楚 & \texttt{arrangement-lead-vocal} \\
双唱 & \texttt{arrangement-vocal-double} \\
和声 & \texttt{arrangement-vocal-harmony} \\
叠唱与装饰句 & \texttt{arrangement-vocal-stack} \\
夜半落地 & \texttt{arrangement-vocal-yeban} \\
常见错误与修正 & \texttt{arrangement-vocal-mistakes} \\
练习 & \texttt{arrangement-vocal-exercise} \\
\bottomrule
\end{longtable}
\end{helpmapping}
